Since your book transitions from these philosophical questions into a rigorous logic section and then into synthetic geometry, this outline is designed to build a "logical ramp." It moves from the abstract nature of a point to the practical mechanics of your axiomatic system.Here is a detailed outline for Part I: Philosophical Foundations.\part{Philosophical Foundations}\chapter{The Ontology of Geometry}Focus: The "What." Defining the status of the objects you are about to manipulate.\section{Existence: The Reality of Mathematical Objects}Platonism vs. Nominalism: Do points $A, B, C$ exist in a non-physical realm, or are they just useful linguistic shorthand?Structuralism: The argument that the relations (incidence, betweenness) are more "real" than the objects themselves.\section{Origin: Discovery and Invention}Is the "Hjelmslev quadrilateral" something you found in the fabric of logic, or a human construction designed to solve a problem?\section{Modality: Necessity, Possibility, and Contingency}Geometric Necessity: Why must the interior angles of a triangle sum to $180^\circ$ in one system but not another?Essential vs. Accidental Properties: Is "having three sides" essential to a triangle, or is it a byproduct of more fundamental axioms?\section{Identity: The Nature of Objects and Relations}Indiscernibility: If two points share all the same properties (incidence), are they the same point? (The Identity of Indiscernibles).\chapter{Epistemology of Mathematics}Focus: The "How we know." Justifying the truth-claims of your theorems.\section{The Nature of Geometric Truth}A Priori Knowledge: Why geometry is the "gold standard" of knowledge that doesn't require physical experiment.\section{The Anatomy of Deduction}Axioms vs. Postulates: Re-evaluating the classical distinction (self-evident truths vs. required assumptions).The Role of Definitions: Are definitions "abbreviations" or do they create new mathematical "entities"?\section{Proof, Justification, and Rigor}The "Gap" Problem: Historical analysis of where Euclid relied on "intuition" and how modern rigor (Hilbert/Pasch) closes those gaps.\section{The Epistemic Status of Diagrams}Visual Evidence: Can a TikZ diagram be "truth-bearing," or is it merely a psychological aid?The "Lettered Diagram": How the interaction between text and image creates a unique form of mathematical understanding.\chapter{Methodology and Metatheory}Focus: The "Rules of the Game." Setting the stage for your specific formal system.\section{The Synthetic Paradigm}Anti-Reductionism: Why we should study geometry on its own terms rather than reducing it to real numbers and $x, y$ coordinates.The Purity of Method: The philosophical value of proving a theorem using the fewest "outside" tools possible.\section{Criteria for Axiomatic Selection}Parsimony (Occam’s Razor): The drive for a "minimal" set of axioms.Aesthetics of Symmetry: Choosing axioms that "feel" balanced and reflect the underlying symmetry of space.\section{Metatheoretical Desiderata}Consistency: The requirement that your system never proves $P$ and $\neg P$.Independence: The beauty of a "lean" system where no axiom is redundant.Categoricity: Does your system describe the one true plane, or a family of possible spaces?\section{The Computational Turn: Formal Verification}The Lean Perspective: How translating geometry into a formal language (like Lean 4) changes our philosophical definition of a "proof."Machine vs. Human: Can a proof be "true" if it is too complex for a human to hold in their mind simultaneously?\chapter{Aesthetics and Teleology}Focus: The "Why." The human element of your monograph.\section{The Elegance of Synthetic Proof}Defining "Mathematical Beauty" through the lens of geometry—symmetry, surprise, and inevitability.\section{Teleology: The Goal of this System}Is the goal to teach, to formalize, or to explore a specific logical niche (like non-Euclidean quadrilaterals)?Pedagogical vs. Research Intent: How the "end goal" shapes the way axioms are ordered.

axiomatic-geometry-foundations/
│
├── main/
│   ├── main.tex                      % MAIN COMPILATION ENTRY
│   ├── draft.tex                    % debug / experimental build
│   └── annotated.tex               % TODO + notes compiled version
│
├── styles/
│   ├── geometry-foundations.sty    % custom macros + environments
│   ├── notation.sty                % math + logic formatting rules
│   ├── theorem-setup.sty           % axiom/def/theorem styling
│   └── layout.sty                  % spacing, headers, structure
│
├── content/
│   │
│   ├── frontmatter/
│   │   ├── title.tex
│   │   ├── abstract.tex
│   │   ├── preface.tex
│   │   ├── notation-guide.tex
│   │   └── roadmap.tex
│   │
│   ├── body/
│   │   │
│   │   ├── philosophy/
│   │   │   ├── ontology.tex
│   │   │   ├── epistemology.tex
│   │   │   └── methodology.tex
│   │   │
│   │   ├── logic/
│   │   │   ├── first-order.tex
│   │   │   ├── syntax.tex
│   │   │   ├── semantics.tex
│   │   │   └── deduction.tex
│   │   │
│   │   ├── formal-system/
│   │   │   ├── signatures.tex
│   │   │   ├── language.tex
│   │   │   ├── well-formedness.tex
│   │   │   └── models.tex
│   │   │
│   │   ├── set-theoretic/
│   │   │   ├── primitives.tex
│   │   │   ├── axioms/
│   │   │   │   ├── equality.tex
│   │   │   │   ├── betweenness.tex
│   │   │   │   ├── congruence.tex
│   │   │   │   └── incidence.tex
│   │   │   ├── constructions/
│   │   │   │   ├── lines.tex
│   │   │   │   ├── segments.tex
│   │   │   │   └── intersections.tex
│   │   │   ├── objects.tex
│   │   │   └── theorems.tex
│   │   │
│   │   ├── type-theoretic/
│   │   │   ├── foundations.tex
│   │   │   ├── identity-types.tex
│   │   │   ├── constructive-geometry.tex
│   │   │   └── computation.tex
│   │   │
│   │   ├── category-theoretic/
│   │   │   ├── categories.tex
│   │   │   ├── morphisms.tex
│   │   │   ├── geometry-as-functor.tex
│   │   │   └── topos-geometry.tex
│   │   │
│   │   └── equivalence/
│   │       ├── translations.tex
│   │       ├── equivalence-theorems.tex
│   │       └── philosophical-implications.tex
│   │
│   ├── backmatter/
│   │   ├── glossary.tex
│   │   ├── notation-index.tex
│   │   ├── bibliography/
│   │   │   └── refs.bib
│   │   ├── index.tex
│   │   └── appendices/
│   │       ├── appendix-a.tex
│   │       ├── appendix-b.tex
│   │       └── proofs.tex
│
├── figures/
│   ├── tikz/
│   ├── diagrams/
│   └── generated/
│
├── code/
│   ├── lean/
│   │   ├── geometry.lean
│   │   ├── axioms.lean
│   │   └── equivalence.lean
│   │
│   ├── python/
│   │   └── visualization.py
│   │
│   └── validation/
│       └── consistency-checks.py
│
├── meta/
│   │
│   ├── README.md
│   ├── project-overview.md
│   ├── roadmap.md
│   ├── todo.md
│   ├── ideas.md
│   ├── notation-spec.md
│   ├── writing-style-guide.md
│   ├── folder-architecture.md
│   ├── compilation-notes.md
│   ├── design-decisions.md
│   └── change-log.md
│
├── build/
│   ├── output.pdf
│   ├── aux/
│   └── logs/
│
├── .vscode/
│   └── axiomatic-geometry.code-workspace
│
├── .gitignore
└── README.md

├── FRONT MATTER
│   ├── Title Page
│   ├── Copyright Page
│   ├── Dedication (optional)
│   ├── Preface
│   ├── Acknowledgments
│   ├── Abstract / Overview
│   └── Table of Contents
│
├── PART I: Philosophical Foundations
│   │
│   ├── Chapter: The Ontology of Geometry
│   │   ├── Section: Existence — The Reality of Mathematical Objects
│   │   ├── Section: Origin — Discovery and Invention
│   │   ├── Section: Modality — Necessity, Possibility, and Contingency
│   │   └── Section: Identity — The Nature of Objects and Relations
│   │
│   └── Chapter: Epistemology of Mathematics
│       ├── Section: The Nature of Truth
│       ├── Section: Axioms, Theorems, and Structure
│       ├── Section: Definitions
│       └── Section: Proof, Justification, and Rigor
│
├── PART II: Logical Foundations
│   │
│   ├── Chapter: First-Order Logic
│   │   ├── Section: Syntax of Formal Languages
│   │   ├── Section: Logical Connectives
│   │   ├── Section: Variables and Quantifiers
│   │   └── Section: Rules of Inference and Deduction
│   │
│   ├── Chapter: The Signature of Geometry
│   │   ├── Section: Formal Signatures
│   │   ├── Section: Geometric Language 𝓛_geo
│   │   ├── Section: Well-Formed Formulas
│   │   └── Section: Typed Extensions of the Signature (Preview)
│   │
│   └── Chapter: The Axiomatic System
│       ├── Section: Primitive Objects
│       ├── Section: Primitive Relations
│       ├── Section: Axioms
│       ├── Section: Derived Notions
│       ├── Section: Derived Truths
│       └── Section: Models and Structures
│
├── PART III: Set-Theoretic Geometry
│   │
│   ├── Chapter: Primitive Structure
│   │   ├── Section: Points
│   │   ├── Section: Relations (Incidence, Equality, Betweenness, Congruence)
│   │   └── Section: The Plane as a Structure
│   │
│   ├── Chapter: Axioms
│   │   ├── Section: Equality Axioms
│   │   ├── Section: Betweenness Axioms
│   │   ├── Section: Congruence Axioms
│   │   ├── Section: Incidence Axioms
│   │   └── Section: Plane Axioms
│   │
│   ├── Chapter: Geometric Constructions
│   │   ├── Section: Line Segments
│   │   ├── Section: Lines
│   │   ├── Section: Rays
│   │   └── Section: Intersection and Extension
│   │
│   ├── Chapter: Geometric Objects
│   │   ├── Section: Triangles
│   │   ├── Section: Polygons
│   │   └── Section: Derived Relations
│   │
│   └── Chapter: Theorems of the System
│       ├── Section: Basic Properties
│       ├── Section: Order Properties
│       └── Section: Congruence Results
│
├── PART IV: Type-Theoretic Geometry
│   │
│   ├── Chapter: Type-Theoretic Primitives
│   │   ├── Section: Points as Types
│   │   ├── Section: Relations as Dependent Types
│   │   └── Section: Identity Types and Equality
│   │
│   ├── Chapter: Constructive Axioms
│   │   ├── Section: Equality as Identity
│   │   ├── Section: Betweenness as a Type
│   │   ├── Section: Constructive Existence
│   │   └── Section: Apartness vs Inequality
│   │
│   ├── Chapter: Geometric Constructions
│   │   ├── Section: Dependent Lines (Lines as Types)
│   │   ├── Section: Constructive Segments
│   │   └── Section: Proof-Relevant Geometry
│   │
│   └── Chapter: Computational Interpretation
│       ├── Section: Proofs as Constructions
│       ├── Section: Geometry as Programs
│       └── Section: Algorithmic Existence
│
├── PART V: Category-Theoretic Geometry
│   │
│   ├── Chapter: Categorical Foundations
│   │   ├── Section: Categories and Morphisms
│   │   ├── Section: Objects of Geometry
│   │   └── Section: Points as Morphisms
│   │
│   ├── Chapter: Geometric Structure in Categories
│   │   ├── Section: Lines as Objects
│   │   ├── Section: Incidence as Pullbacks
│   │   └── Section: Congruence as Isomorphism
│   │
│   ├── Chapter: Advanced Structure
│   │   ├── Section: Groupoids of Transformations
│   │   ├── Section: Symmetry and Automorphisms
│   │   └── Section: Functorial Geometry
│   │
│   └── Chapter: Topos-Theoretic Interpretation
│       ├── Section: Geometry in a Topos
│       ├── Section: Internal Logic
│       └── Section: Synthetic Geometry
│
├── PART VI: Equivalence and Interpretation
│   │
│   ├── Chapter: Translation Between Foundations
│   │   ├── Section: Set-Theoretic vs Type-Theoretic Interpretation
│   │   ├── Section: Relations vs Types
│   │   └── Section: Truth vs Inhabitance
│   │
│   ├── Chapter: Equivalence of Geometries
│   │   ├── Section: Preservation of Theorems
│   │   ├── Section: Structural Correspondence
│   │   └── Section: Limits of Equivalence
│   │
│   └── Chapter: Mathematics as Invention
│       ├── Section: Multiple Realizations of Structure
│       ├── Section: Foundations as Encoding Choices
│       └── Section: Against Uniqueness of Mathematical Reality
│
├── BACK MATTER
│   ├── Appendix A: Formal Proofs / Extended Constructions
│   ├── Appendix B: Lean / Formalization Code (optional)
│   ├── Appendix C: Additional Axiomatic Variants (optional)
│   ├── Bibliography
│   ├── Index
│   └── Notation Index (optional but VERY useful)
│
└── END
\end{document}

% ==================================================
\chapter{Project Documentation and Style Specification}

This chapter defines the formal conventions, structural rules, and development guidelines for the construction of this axiomatic geometry system. These rules govern not only the mathematical content, but also the organization of files, notation, and LaTeX usage across the project.

% ==================================================
\section{Project Philosophy}

This project is structured as a multi-foundational reconstruction of geometry. The goal is not only to formalize Euclidean-style geometry, but to express the same underlying mathematical structure across multiple foundational frameworks:

\begin{itemize}
    \item Set-theoretic formulation (classical foundation)
    \item Type-theoretic formulation (constructive / computational interpretation)
    \item Category-theoretic formulation (structural / relational abstraction)
\end{itemize}

The guiding principle is:

\begin{quote}
Mathematics is treated as a family of equivalent structural encodings rather than a single foundational ontology.
\end{quote}

% ==================================================
\section{Project Folder Architecture}

The project is organized into a modular hierarchy reflecting conceptual dependency:

\begin{verbatim}
axiomatic-geometry-foundations/
│
├── main/                  Main compilation entry point
├── styles/               Custom LaTeX macros and formatting system
├── content/              Mathematical content (core system)
│   ├── frontmatter/
│   ├── body/
│   │   ├── philosophy/
│   │   ├── logic/
│   │   ├── set-theoretic/
│   │   ├── type-theoretic/
│   │   └── category-theoretic/
│   └── backmatter/
│
├── figures/              TikZ and diagram assets
├── code/                 Formalization / Lean / computation experiments
├── meta/                 Documentation, notes, TODOs, design specs
└── workspace/            IDE configuration (.code-workspace)
\end{verbatim}

\subsection{Meta Directory}

The \texttt{meta/} directory contains non-mathematical but structurally essential information:

\begin{itemize}
    \item Style guide (this document)
    \item TODO lists and development roadmap
    \item Design notes and philosophical decisions
    \item Experimental ideas and rejected approaches
    \item Versioning and changelogs
\end{itemize}

% ==================================================
\section{LaTeX Package Architecture}

The project uses a layered LaTeX system.

\subsection{Core Packages}

\begin{itemize}
    \item \texttt{amsmath, mathtools} — advanced mathematical typesetting
    \item \texttt{amsthm} — theorem and axiom environments
    \item \texttt{hyperref} — cross-referencing and navigation
    \item \texttt{comment} — modular compilation control
\end{itemize}

\subsection{Custom Style Layer (\texttt{styles/})}

A dedicated style file (\texttt{geometry-foundations.sty}) is used to define:

\begin{itemize}
    \item Global notation rules
    \item Custom macros for relations (\texttt{\textbackslash Bet}, \texttt{\textbackslash Cong})
    \item Formatting rules for axioms and definitions
    \item Consistent spacing behavior in math mode
\end{itemize}

\subsection{Design Principle}

All semantic structure should be separated from presentation:

\begin{quote}
Mathematical meaning is encoded in macros; typography is delegated to style files.
\end{quote}

% ==================================================
\section{Notation and Typesetting Rules}

\subsection{Display Mathematics}

\begin{itemize}
    \item Always use:
\end{itemize}

\begin{verbatim}
\[
    a + b = c
\]
\end{verbatim}

\begin{itemize}
    \item Never use inline \texttt{\$\$...\$\$}
    \item Always include spaces inside expressions:
\end{itemize}

\[
 a + b = c
\]

\subsection{Set Notation}

Use scalable delimiters:

\[
\left\{ x \mid x > 0 \right\}
\]

Never use fixed-size braces for non-trivial expressions.

\subsection{Relations}

Standard notation:

\begin{itemize}
    \item Equality: $A = B$
    \item Betweenness: $\mathcal{B}(A,B,C)$
    \item Congruence: $A \cong B$
\end{itemize}

All relations must be treated as first-class logical objects.

% ==================================================
\section{Structural Formatting Rules}

\subsection{Section Separation}

\begin{itemize}
    \item One blank line before and after every \texttt{\textbackslash section}
    \item Use comment dividers for major structural blocks:
\end{itemize}

\begin{verbatim}
% =========================
\section{Example}
% =========================
\end{verbatim}

\subsection{Indentation}

\begin{itemize}
    \item Environments must be indented
    \item Nested logical structures must preserve hierarchy visually
\end{itemize}

\subsection{Itemization Rules}

\begin{itemize}
    \item One item per line in \texttt{itemize} and \texttt{enumerate}
    \item No inline multi-item lists
\end{itemize}

% ==================================================
\section{Axiom and Definition Conventions}

\subsection{Axioms}

\begin{itemize}
    \item Always use \texttt{axiom} environment
    \item Every axiom must be globally self-contained
    \item No external justification inside axiom blocks
\end{itemize}

\subsection{Definitions}

\begin{itemize}
    \item Definitions must introduce exactly one new concept
    \item Must include formal symbolic expression when possible
\end{itemize}

\subsection{Naming}

\begin{itemize}
    \item Axioms must have descriptive titles:
    \item Example: \textit{Symmetry of Equality}
\end{itemize}

% ==================================================
\section{Commenting and Meta-Annotating}

Comments are not decorative; they are structural metadata.

\subsection{Allowed Uses}

\begin{itemize}
    \item External references (papers, Wikipedia, ideas)
    \item Alternative formulations
    \item Proof sketches
    \item Future expansions
\end{itemize}

\subsection{Forbidden Uses}

\begin{itemize}
    \item Redundant restatement of visible content
    \item Informal explanations of obvious steps
\end{itemize}

% ==================================================
\section{Code and Formalization Layer}

The \texttt{code/} directory is reserved for formal systems such as:

\begin{itemize}
    \item Lean 4 formal proofs
    \item Type-theoretic encodings
    \item Computational geometry simulations
\end{itemize}

Principle:

\begin{quote}
Every informal definition should eventually have a formal counterpart.
\end{quote}

% ==================================================
\section{TODO System and Development Workflow}

All incomplete or speculative content must be explicitly tracked.

\subsection{TODO Categories}

\begin{itemize}
    \item Structural TODO — missing definitions or axioms
    \item Logical TODO — incomplete proofs or derivations
    \item Design TODO — unclear notation or ambiguity
    \item Expansion TODO — future theoretical extensions
\end{itemize}

\subsection{Format}

\begin{verbatim}
% TODO: clarify definition of congruence class
% IDEA: reinterpret betweenness categorically
\end{verbatim}

\subsection{Rule}

No unresolved TODO may remain in finalized chapters.

% ==================================================
\section{Cross-Foundational Consistency Rule}

All three foundational systems must satisfy:

\begin{quote}
They must encode identical geometric content under different semantic interpretations.
\end{quote}

This implies:

\begin{itemize}
    \item Set theory: relational membership structure
    \item Type theory: constructive witness structure
    \item Category theory: morphism-preserving structure
\end{itemize}

% ==================================================
\section{Versioning and Evolution}

The system is expected to evolve.

Recommended versioning:

\begin{itemize}
    \item v0.x — structural design phase
    \item v1.x — complete set-theoretic system
    \item v2.x — cross-foundational equivalence established
\end{itemize}

% ==================================================
\section{Design Extensions (Future Work)}

Potential expansions include:

\begin{itemize}
    \item Formal category-theoretic semantics of betweenness
    \item Computational verification of axioms
    \item Homotopy-theoretic interpretation of geometric equality
    \item Synthetic geometry via topos theory
\end{itemize}

% ==================================================
\section{Closing Principle}

This project is not only a mathematical construction, but a study of representation itself:

\begin{quote}
The same geometric structure may exist simultaneously as set-theoretic relations, type-theoretic constructions, and categorical morphisms.
\end{quote}